% =====================================================================
% Figure contract
% 核心结论：CUA 是 Tool-use agent 的子类，并以 GUI 主通道边界区别于 RPA 与纯 CLI/API agent。
% 图原型：中等复杂度、纯结构的嵌套集合图；用 fit 构造 Tool-use agent > CUA 两层包含关系。
% 证据层级：hero = 主色重描边的 CUA 框与四个平台卡；支撑 = GUI 异名、RPA 前身、CLI 姊妹范式。
% 能不能更少？：不能；每个框和三条关系线均对应 §2.2--2.3 的一个定义性关系。
%
% Step ① 设计文档（Form A / 中等档）
% 领域 / 格式：Computer-Use Agents 中文综述；TikZ standalone + XeLaTeX。
% 目标印刷宽度 W_print = 19.0 cm；目标印刷最小字号 = 7.0 pt；设计最小字号 = 7.6 pt。
% 布局：标题在上；RPA 在左；Tool-use 浅灰父框居中，内含 CUA、GUI Agent 与 CLI/coding agent；
%       CUA 内四卡水平并排；底部为整宽定义性边界注释条。目标画布宽度 < 20 cm。
% 连线：RPA -> CUA 为实线短箭头；CUA <-> GUI 为虚线双向；Hybrid -- CLI 为垂直实线无箭头。
% 空间：三条连线各占独立走廊；标签放在线旁空白处；跨框只穿越与关系有关的父/子边界。
% 强调：CUA 用主色 1.4 pt 描边；Tool-use 用浅灰底；相邻概念卡降低饱和度与描边重量。
% ASCII 草图：
%             CUA 与相邻概念的关系（§2.2–2.3）
%   [RPA] --> + Tool-use agent --------------------------------------+
%              | + CUA -----------------------+ <--> [GUI Agent 范式] |
%              | | Web | Mobile | OS | Hybrid |                      |
%              | +-----------------------------+                      |
%              |                         | 通道互补                   |
%              |                  [CLI / coding agent]                |
%              +-----------------------------------------------------+
%   [定义性边界注释条：GUI 主通道边界]
% Pre-flight：同行四卡由 positioning 构造；两层容器由 fit 构造；全部端点使用 node anchor；
%             无旋转中文；标签与连线错位放置；内层卡到 CUA 边界留足 inner sep。
% =====================================================================

\documentclass[border=8pt]{standalone}
\usepackage{fontspec}
\usepackage{xeCJK}
\usepackage{xcolor}
\usepackage{tikz}
\setCJKmainfont{Heiti SC}
\setCJKsansfont{Heiti SC}
\usetikzlibrary{arrows.meta,bending,positioning,calc,backgrounds,fit}

% Low-saturation academic palette, shared with cua-survey-fig1-overview.tex.
\definecolor{ink}{HTML}{27364A}
\definecolor{muted}{HTML}{748092}
\definecolor{mainLine}{HTML}{4F667C}
\definecolor{mainFill}{HTML}{EAF0F4}
\definecolor{heroFill}{HTML}{F1F3F6}
\definecolor{heroBorder}{HTML}{3F566C}
\definecolor{blueLine}{HTML}{6D879E}
\definecolor{blueFill}{HTML}{E7EEF4}
\definecolor{purpleLine}{HTML}{817594}
\definecolor{purpleFill}{HTML}{EEEAF2}
\definecolor{goldLine}{HTML}{9B875B}
\definecolor{goldFill}{HTML}{F2EEE3}
\definecolor{tealLine}{HTML}{648D8A}
\definecolor{tealFill}{HTML}{E5F0EE}

\pgfdeclarelayer{background}
\pgfsetlayers{background,main}

\tikzset{
  every node/.style={
    font=\sffamily\fontsize{8.4}{10.2}\selectfont,
    text=ink
  },
  platform/.style={
    rectangle,
    rounded corners=4pt,
    line width=0.7pt,
    minimum width=1.875cm,
    minimum height=1.15cm,
    text width=1.62cm,
    align=center,
    inner sep=5pt,
    font=\sffamily\fontsize{8.1}{9.6}\selectfont\bfseries
  },
  sidecard/.style={
    rectangle,
    rounded corners=5pt,
    line width=0.7pt,
    align=center,
    inner sep=7pt
  },
  relationtag/.style={
    font=\sffamily\fontsize{7.4}{8.8}\selectfont,
    text=muted,
    inner sep=1.5pt
  },
  arrow short/.style={
    -{Stealth[length=3pt 1.0,width'=0pt 0.5,sep=0pt 0.3]},
    line width=0.8pt,
    shorten >=1pt,
    shorten <=0pt,
    color=mainLine
  },
  bidirectional/.style={
    arrows={Stealth[length=3.2pt 1.0,width'=0pt 0.55]-Stealth[length=3.2pt 1.0,width'=0pt 0.55]},
    densely dashed,
    line width=0.85pt,
    shorten >=2pt,
    shorten <=2pt,
    color=purpleLine
  },
  complement/.style={
    line width=0.85pt,
    color=blueLine,
    shorten >=2pt,
    shorten <=2pt
  }
}

\begin{document}
\sffamily
\begin{tikzpicture}

% CUA platform cards: positioned first so both containing frames can be fit by construction.
\node[platform,fill=blueFill,draw=blueLine] (web) {Web Agent};
\node[platform,fill=tealFill,draw=tealLine,right=0.22cm of web] (mobile) {Mobile Agent};
\node[platform,fill=goldFill,draw=goldLine,right=0.22cm of mobile] (os) {OS Agent};
\node[platform,fill=purpleFill,draw=purpleLine,right=0.22cm of os] (hybrid) {GUI+API\\Hybrid};

\node[
  anchor=south west,
  font=\sffamily\fontsize{9.3}{11.0}\selectfont\bfseries,
  text=heroBorder
] (cuaTitle) at ([yshift=0.42cm]web.north west) {CUA — 跨平台统称};

% Invisible bounds nodes make the nested geometry available before background painting.
\node[fit=(cuaTitle)(web)(mobile)(os)(hybrid),inner xsep=0.40cm,inner ysep=0.34cm] (cuaBounds) {};

\node[
  sidecard,
  fill=purpleFill,
  draw=purpleLine,
  minimum width=2.45cm,
  minimum height=1.10cm,
  right=1.65cm of cuaBounds,
  font=\sffamily\fontsize{8.5}{10.0}\selectfont\bfseries
] (gui) {GUI Agent 范式};

\node[
  sidecard,
  fill=blueFill,
  draw=blueLine,
  minimum width=5.10cm,
  minimum height=1.03cm,
  below=1.08cm of hybrid,
  font=\sffamily\fontsize{8.0}{9.6}\selectfont
] (cli) {CLI / coding agent — 无视觉输入，\\姊妹范式};

\node[
  anchor=south west,
  font=\sffamily\fontsize{8.8}{10.5}\selectfont\bfseries,
  text=heroBorder
] (toolTitle) at ([yshift=0.78cm]cuaBounds.north west)
  {Tool-use agent — 任意 function-calling / API 调用 agent（父类）};

\node[fit=(toolTitle)(cuaBounds)(gui)(cli),inner xsep=0.45cm,inner ysep=0.38cm] (toolBounds) {};

\node[
  sidecard,
  anchor=east,
  fill=heroFill,
  draw=muted!78,
  minimum width=3.30cm,
  minimum height=1.52cm,
  font=\sffamily\fontsize{7.8}{9.4}\selectfont
] (rpa) at ([xshift=-1.40cm]cuaBounds.west)
  {RPA / Programming-by-\\Demonstration — 录制重放，\\无语义泛化};

% The two nested containers are painted behind all foreground text and cards.
\begin{pgfonlayer}{background}
  \node[
    rectangle,
    rounded corners=8pt,
    draw=muted!65,
    fill=heroFill,
    line width=0.8pt,
    fit=(toolTitle)(cuaBounds)(gui)(cli),
    inner xsep=0.45cm,
    inner ysep=0.38cm
  ] {};
  \node[
    rectangle,
    rounded corners=7pt,
    draw=heroBorder,
    fill=mainFill,
    line width=1.4pt,
    fit=(cuaTitle)(web)(mobile)(os)(hybrid),
    inner xsep=0.40cm,
    inner ysep=0.34cm
  ] {};
\end{pgfonlayer}

% Relationships. Each line occupies its own corridor and crosses only a relevant boundary.
\draw[arrow short]
  (rpa.east) -- (cuaBounds.west)
  node[relationtag,pos=0.31,above=4pt,font=\sffamily\fontsize{7.0}{8.4}\selectfont] {历史前身};

\draw[bidirectional]
  (cuaBounds.east) -- (gui.west)
  node[relationtag,midway,above=4pt,text=purpleLine] {同一范式异名};

\draw[complement]
  (hybrid.south) -- (cli.north)
  node[relationtag,midway,right=5pt,text=blueLine] {通道互补};

% Title and full-width definition bar.
\coordinate (fullCenter) at ($(rpa.west)!0.5!(toolBounds.east)$);
\node[
  anchor=south,
  font=\sffamily\fontsize{12.0}{14.2}\selectfont\bfseries,
  text=heroBorder
] at ($(fullCenter |- toolBounds.north)+(0,0.70cm)$)
  {CUA 与相邻概念的关系（§2.2–2.3）};

\node[
  rectangle,
  rounded corners=4pt,
  draw=muted!55,
  fill=mainFill!55,
  line width=0.6pt,
  minimum width=18.45cm,
  minimum height=1.15cm,
  align=center,
  inner sep=6pt,
  anchor=north,
  font=\sffamily\fontsize{7.8}{9.6}\selectfont\bfseries
] at ($(fullCenter |- toolBounds.south)+(0,-0.55cm)$)
  {定义性边界：GUI 是否为主要观察与操作通道（§2.1、§2.3）\\[-1pt]
   ——把 CUA 与纯 API/CLI agent 和 RPA 区分开};

\end{tikzpicture}
\end{document}
